\begin{pastexam}[label={px:ntu-113-dsa-02}]{臺大資工所 113 學年度 第 2 題}
（複選；每個選項單獨計分）

(5 points) A Maximum Binary Tree (abbreviated as MaxBT) is defined as a binary tree that adheres to
the max-heap property without the necessity of following the complete binary tree property. In a
MaxBT, each node $p$ consists of a value ($p$.data) and two links ($p$.left and $p$.right) pointing to
the left and right children, respectively. A link to nil indicates the absence of further children.

The provided pseudocode, although incomplete, implements a recursive function for merging two MaxBTs
into a single MaxBT. This function takes the root nodes of the two MaxBTs ($\mathit{root1}$ and
$\mathit{root2}$) as input and returns the root node of the merged MaxBT.
\begin{center}
\begin{tabular}{rl}
1: & function MergeMaxBinaryTrees($\mathit{root1}$, $\mathit{root2}$) \\
2: & \quad if $\mathit{root1}$ = nil then return \blank{a} end if \\
3: & \quad if $\mathit{root2}$ = nil then return \blank{b} end if \\
4: & \quad if $\mathit{root1}$.data $>$ $\mathit{root2}$.data then \\
5: & \quad\quad \blank{c} $\leftarrow$ MergeMaxBinaryTrees(\blank{d}, \blank{e}) \\
6: & \quad\quad return \blank{f} \\
7: & \quad else \\
8: & \quad\quad \blank{g} $\leftarrow$ MergeMaxBinaryTrees(\blank{h}, \blank{i}) \\
9: & \quad\quad return \blank{j} \\
10: & \quad end if \\
11: & end function
\end{tabular}
\end{center}
Assuming all the blanks in the above function have been appropriately filled to ensure its proper
functioning, please select the correct description(s) below.
\begin{enumerate}[label=\Alph*.]
\item (b) and (f) are the same.
\item (c), (d), and (e) are all distinct.
\item In the cases of (c) and (g), each has only one correct answer.
\item Exchanging the filled content of (h) and (i) leads to a correct implementation.
\item The function works correctly even when the two arguments passed into it refer to the same
tree.
\end{enumerate}
\end{pastexam}

\begin{pastsoln}
答案：A、D。

先把函式補完。 兩棵樹都不空時，根比較大的那個當新的根（$>$ 不成立時選 $\mathit{root2}$，
相等也可以），把它的某一棵子樹和另一整棵樹遞迴合併，結果接回同一邊。 這和資結書
虛擬碼~8.9 左傾樹的 \proc{Meld} 是同一個架構，只是少了比較 $s$、對調左右那兩步。
一種正確的填法：

\begin{pseudo}[unbreakable]{一種正確的填法}
\Function{MergeMaxBinaryTrees}{$\mathit{root1}, \mathit{root2}$}
    \If{$\mathit{root1} = \NIL$} \Return $\mathit{root2}$ \Comment{(a)} \EndIf
    \If{$\mathit{root2} = \NIL$} \Return $\mathit{root1}$ \Comment{(b)} \EndIf
    \If{$\mathit{root1}.\mathit{data} > \mathit{root2}.\mathit{data}$}
        \State $\mathit{root1}.\mathit{right} \gets$ \Call{MergeMaxBinaryTrees}{$\mathit{root1}.\mathit{right}, \mathit{root2}$} \Comment{(c)(d)(e)}
        \State \Return $\mathit{root1}$ \Comment{(f)}
    \Else
        \State $\mathit{root2}.\mathit{right} \gets$ \Call{MergeMaxBinaryTrees}{$\mathit{root1}, \mathit{root2}.\mathit{right}$} \Comment{(g)(h)(i)}
        \State \Return $\mathit{root2}$ \Comment{(j)}
    \EndIf
\EndFunction
\end{pseudo}

正確性對兩棵樹的節點總數歸納：遞迴呼叫的兩棵樹節點總數少了一個（新的根不在裡面），由歸納假
設回傳一棵 MaxBT，它的節點都來自原本的兩棵樹，所以都不大於新的根；接回去以後新的根仍然不小
於兩棵子樹的所有節點，整棵是 MaxBT。

逐個選項：
\begin{itemize}
\item[A.] 對。 (b) 是 $\mathit{root2}$ 為空時回傳 $\mathit{root1}$；(f) 是 $\mathit{root1}$ 當新的
根時回傳 $\mathit{root1}$。 兩格都必須是 $\mathit{root1}$。
\item[B.] 錯。 (c) 是被改寫的那個欄位，它原本指的子樹必須當作遞迴的一個參數傳進去，否則那棵子
樹就從樹上掉了。 所以 (c) 一定和 (d)、(e) 其中一個相同（例如上面的 (c) 和 (d) 都是
$\mathit{root1}$.right）。
\item[C.] 錯。 MaxBT 不要求形狀，(c) 填 $\mathit{root1}$.left（同時 (d) 或 (e) 改成
$\mathit{root1}$.left）也一樣正確；(g) 同理。 每一格都至少有兩種正確答案。
\item[D.] 對。 上面的正確性論證沒有用到兩個參數的順序：交換 (h)、(i) 之後，遞迴呼叫仍然是合併同
樣兩棵樹，回傳的仍然是一棵包含它們全部節點的 MaxBT。 相等的情況也沒問題，相等時選誰當根都
滿足最大堆積性質。
\item[E.] 錯。 兩個參數是同一棵樹時，資料相等，走 else，把 $\mathit{root2}$.right 改成「這棵樹和它
自己的右子樹合併」的結果，同一批節點被接到兩個地方。 最簡單的例子是只有一個節點 $r$：
\proc{MergeMaxBinaryTrees}$(r, r)$ 走 else，$r.\mathit{right} \gets$
\proc{MergeMaxBinaryTrees}$(r, \NIL) = r$，$r$ 的右子節點變成自己，不再是一棵樹。
\end{itemize}

補充：沒有像左傾樹那樣控制右邊路徑的長度，這個合併最壞要沿著很長的路徑遞迴，時間是
$\Theta(n)$，而左傾樹保證 $O(\lg n)$（資結書定理~8.2）。
\end{pastsoln}
